\section{The Man Who Knew Too Much}

\begin{quotex}
One might say that set theory is to mathematics what quantum theory is to physics; it is what remains when every last vestige of ``being'' has been exorcised.
\flright{\textsc{Wolfgang Smith}}
\end{quotex}


It is hardly uncommon to listen to amateurish metaphysicians on the Internet who are fixated on quantum theory and set theory. One fellow droned on and on about an ``actual infinity'', whatever that might mean. There was no way to interrupt, so I had to click pause. Obviously, such a concept is incoherent because there is only one Infinity. He is just one of many with a new Theory of Everything. Often they are merely the sequel to \textbf{Alfred Hitchcock}'s \emph{The Man Who Knew Too Much}: The Man Who Understood Too Little.

What Wolfgang Smith means is that those concepts are so abstract that they become essentially meaningless. Quantum particles do not ``exist'' as things; rather they are described by possible quantum states. Only when ``measured'', what Dr. Smith calls vertical causation, do they become beings. Similarly, as set is a mathematical concept with no properties other than its members.

It is the height of arrogance, inspired by presentism, to presume that no one in the history of the human race ever understood anything about the world, the human race, and God until quantum and set theories were invented, yes, ``invented'', not ``discovered''.

To know God, it is best to rely on the esoterists, saints, prophets, etc., who actually knew God, not some innovator with a novel theory about God.

\paragraph{We Have No Bananas Today}


An underemployed philosopher was forced to take a job as produce manager in a German owned supermarket in a small town in Arkansas. Since bananas were sold by the piece rather than by weight, he implemented a system of using sets of bananas in order to track them. He figured that when new bananas arrived, he could just take the union of two sets of bananas to get a count. However, a set must have unique elements but all the bananas looked alike. So he had the stockboy number the bananas to give them an identity.

One day a tribe of monkeys bought all the bananas. When the philosopher asked for inventory, the stockboy answered, ``There are no bananas.''

``You mean,'' explained the philosopher, ``that the bananas are a null set.''

``Yes,'' replied the stockboy, ``we have no bananas today.''

\paragraph{The Consultant}


The philosopher was sacked and is now an Uber driver in the Ozarks. Management brought in a consultant to convert the sets of bananas to their own inventory system. He had to explain the problem as best he could, and here are some passages from his white paper.

A ``set'' of bananas tells you nothing about bananas, since a set is a mathematical construct. As such, it is intermediate between the pure realm of ideas and the sensible realm represented by the bananas. Since the bananas are numbered, as set of 12 bananas is isomorphic to the set of numbers from 1 to 12, for example. Hence, we can dispense with the set of bananas since the set of numbers tells us everything we need to know.

The claim that the universe itself is a set, is no different. At any moment, there are a finite number of physical things in the universe, so the set of things is likewise isomorphic to a finite set of natural numbers, no matter how large it may be. Obviously, calling the universe a set may be tautologically true, but it reveals nothing about the universe.

Moreover, there are two rather vexing problems. The first is how to define a ``thing''. Is a racoon a thing? How about the head of a racoon? This involves questions of mereology, which cannot be solved with set theory alone.

The second problem is that objects are coming into and going out of existence. A set is definitively defined, and its members cannot change. Hence, the universe cannot be a single set, but rather a series of sets.

\paragraph{Epistemic Closure}


By epistemic closure, I mean the notion that the world is self-explanatory. In practice it means the attempt to understand the world from a small number of axioms and any tautologies that necessarily follow from them.

This is clearly false in practice, since there are many things that we don't know about the world. Moreover, if we include mathematics, then there will always be truths that cannot be contained in any system. It is impossible not to notice that all important problems are not resolvable within the universe. Thought is dominated by incessant and futile arguments with no apparent end.

Hence, in principle, the world is incomplete and requires a revelation from outside the world. As we recently showed in the article In The Beginning\footnote{\url{https://gornahoor.net/?p=16477}}, the most important issues about living in the world cannot be solved by science or logic; rather, they must be revealed from outside the world process.

\paragraph{Love is not an Emotion}


I recently watched three men of above average intelligence try to describe love as if it were an emotion. That is so far from the truth of the matter. God is Love, and God is not an emotion. We are commanded to love, and emotions cannot be commanded.

The ultimate point of knowing the world is not knowledge for its own sake, but rather to know the Good. For Plato, it was sufficient to contemplate the Good. However, is the Will is fundamental, then that is insufficient; rather, one must Will the Good. That is what is meant by Love: to Will the Good.

Many people get stuck in a rut of theodicy by misunderstanding physical pain. That is understandable, but spiritual good is unmeasurably more important than material good. Hence, to will the good for someone means to help him spiritually as well as materially.

\paragraph{Divine Simplicity}


God is simple, which means He has no parts. Hence, it is loose talk to claim that a human being is a ``part of God''; that is simply a logical absurdity.

A fortiori, it means that God cannot be the set of all sets. First of all, it is an incoherent notion. A recent comment showed a justification for the claim, but it involved two incompatible notions of containment. The explanation sounded more like the Sokal hoax\footnote{\url{https://en.wikipedia.org/wiki/Sokal\_affair}} than a serious definition.

We cannot worship the set of all sets, nor can it reveal itself. The claim that everything knowable about God is among those sets, so that there is nothing to reveal.

Moreover, as we pointed out, a set is beneath the world of ideas. God, however, is transcendent to everything in the world, hence cannot be a set of any kind.

\paragraph{Concluding Unscientific Postscript}


That is enough for now. Someone's alleged IQ has no bearing on the quality of his system. Even a mountebank like Jordan Peterson considers it necessary to reveal his IQ.

If a system can convince some yobs to believe in God, accept moral responsibility for their postmortem state, and see the influence of Satan in the world, there is some good in it. However, at some point they should tire of feeding on Froot Loops and seek out the solid and nutritious food of Traditional metaphysics.

\url{https://www.youtube.com/watch?v=HygopC4S5W0}

\flrightit{Posted on 2022-08-18 by Cologero}

\begin{center}* * *\end{center}

\begin{footnotesize}
\begin{sffamily}
\texttt{Arthur Konrad on 2022-08-19 at 05:33 said:}

What if our universe is just an elaborate forgery?

\hfill

\texttt{Cologero on 2022-08-19 at 15:20 said:}

Forgery of what? By whom? For what purpose?

\hfill

\texttt{Dimitri on 2022-08-19 at 20:33 said:}

Perhaps he's alluding to the cosmologies of old commonly placed under the umbrella term of ``Gnosticism''? Incidentally, that sort of thinking or at least variations of it seem to be having a resurgence these days, albeit in materialistic terms (e.g. ``the universe is a computer simulation'').

\hfill

\texttt{Michael DeJoy on 2022-08-22 at 09:53 said:}

Cologero, currently there does not appear to be a solution to the Wests' slow motion collapse. Numerous podcasts, blogs, etc. try to argue their points on what needs to be done to stem the tide, however none of these things seem to work. I have been considering the following as a fantastical solution which lies within our grasp. I believe the Catholic Church still has much power and influence in the world that she has let go dormant and often is at odds with its' own original goals. Imagine if one morning Pope Francis woke up and said '' gentlemen, we have been wrong about promoting liberal ideas and we are returning to the ways of our founder. I don't care what governments say, I don't care what news organizations say, I don't care what liberals of all stripes say, we are changing course. Those who agree can stand with us, those against can go their own way.'' I think this would cause a shakeout and restoration of the Church. It would be stronger and rock the world.

\hfill

\texttt{Cologero on 2022-08-22 at 10:59 said:}

What you call a ``slow motion collapse'', most other people call ``progress and freedom''.

If you join the priesthood and are ambitious enough to become pope, then you can implement your program. Otherwise, the pope does not give a rat's ass about your opinion. In any event, the purpose of the Church is to save souls, not to implement some political program.

If you recall, Adam disobeyed God's commandment; since then, many others have followed Adam's example. So if people are quite willing to disobey God, why do you think that they will suddenly obey the pope in your scenario?

There is a Law of Harmony. Everything is in its right place for a reason. This is the world as it is. As I wrote a year ago:

``This is a difficult path; after all, only 3 Knights found the Holy Grail out of the 150 who started. This is meant for those three. Failure comes from this misunderstanding: Tradition is a way of life, not a system of thought.''

While everyone else is looking for a guru with a great system, a few hearty souls go out in search of the Grail without a guidebook. I am looking for Knights, Michael, not whiners.


\hfill

\end{sffamily}
\end{footnotesize}

